\begin{proof}[Proof of Theorem~\ref{main}] We will start with an arbitrary lattice $L$. At each step we will modify the lattice $L$ until we find one which satisfies the equality. 

First note that $n$ vectors $w_1, \dots, w_n$ will generate $L$ if and only if their images $\tilde{w_1}, \dots, \tilde{w_n}$ in $L/tL$ form a basis. 

For each $L_i$, let $W_i \subset L$ be the set of tight generators for $L$ with respect to $L_i$. Then let $\tilde{W}_i \subset L/tL$ be the projection of this set to $L/tL$. It is easy to verify that $\tilde{W}_i$ is a vector subspace of $L/tL$. We wish to choose one vector $\tilde{w_i}$ from each $\tilde{W}_i$ such that the $\tilde{w_1}, \dots, \tilde{w_n}$ form a basis of $L/tL$. In other words, for the subspaces $\tilde{W_i}$, we wish to find a \emph{system of linearly independent representatives}.

We are now equipped to explain the basic idea of the algorithm. Recall that for any $L$, we have
\[
\val(\det(L)) + \sum_{i=1}^n c(L,L_i) \leq A(L_1, L_2, \dots, L_n),
\]
where equality holds exactly when we have a system of linearly independent representatives of the subspaces $\tilde{W_i}$. At each stage, we will either make the left hand side bigger by some integer value, or we will increase the size of the system of linearly independent representatives. The only way this terminates is that we have a system of $n$ linearly independent representatives.

Here is the algorithm for finding $L$. Start with an arbitrary lattice $L$.

We have the subspace $\tilde{W}_i \subset L/tL$ for each $i \leq n$. Start by finding the maximum number of linearly independent representatives $\tilde{w_i}$ for $i \in J$, where $J \subset [n]$. Then we have that by Theorem ~\ref{lin2}, there exists a set $I \subset [n]$ such that 
\[
|J| = \dim \left( \sum_{i\in I} \tilde{W}_i \right) + n - |I|.
\]
Moreover, we have that among the $\tilde{w_i}$, some form a basis for $\tilde{W} \coloneqq  \sum_{i \in I} \tilde{W}_i$, while the rest have indices in $[n] \setminus I$. Thus we may write $J = J_1 \coprod J_2$, where $\tilde{w_i}$ for $i \in J_1$ give a basis of $\tilde{W}$, and $J_2 \subset [n] \setminus I.$

Choose any lift of $\tilde{W}$ to an $\cO$-submodule of $L$. Call this lift $W \subset L$. Now let $L' = t^{-1}W + L$. We claim that
\begin{equation} \label{step}
\val(\det(L)) + \sum_{i=1}^n c(L,L_i) \leq \val(\det(L')) + \sum_{i=1}^n c(L',L_i) .
\end{equation}
To see this, first note that 
\[
\val(\det(L')) = \val(\det(L)) - \dim W.
\]
Also note that 
\begin{align*}
c(L',L_i) &= c(L,L_i) + 1 \text{ for } i \in I
\\
c(L',L_i) &= c(L,L_i) \text{ for } i \in [n] \setminus I.
\end{align*}
Note also that $\dim W \leq |I|$ with equality only when we have a complete system of linearly independent representatives. This yields the inequality~\eqref{step}.

This inequality tells us that, at each stage, if we have not found a complete system of linearly independent representatives, we may modify $L$ to get a lattice $L'$ for which the inequality is strictly closer. We can then iterate this process--find another maximum set of linearly independent representatives; if they still do not span, we can modify the lattice again. The process cannot continue indefinitely, and it will end with a lattice for which we have $n$ linearly independent representatives, and this lattice will satisfy the equality of our theorem.
\end{proof}
